
%%%%%%%%%%  scan idx 46  |  folio 39  %%%%%%%%%%


S G A 7


\underline{E X P O S E  XI}


\underline{COHOMOLOGIE DES INTERSECTIONS COMPLETES}


\underline{par P. DELIGNE}


\underline{1. Cohomologie de l'espace projectif et des intersections compl\`etes}

Soit $X$ une vari\'et\'e projective non singuli\`ere complexe, purement de dimension $n$, et soit $\underline{O}(1)$ un faisceau inversible ample sur $X$. D'apr\`es Akizuki et Nakano [1], il r\'esulte du vanishing theorem de Kodaira [4] que

\begin{equation}\tag{*}
H^i(X,\Omega^j_X (-1)) = 0 \quad \mathrm{pour}\ i+j < n \ .
\end{equation}
D'apr\`es [1] toujours, ce r\'esultat implique le ``th\'eor\`eme de Lefschetz faible'', en cohomologie complexe. On se propose de montrer, par force brutale, que si $X$ est un espace projectif sur un corps de caract\'eristique quelconque, alors (*) reste vrai. Ceci permettra, \`a l'imitation de [1], de comparer les cohomologies de Hodge et de De Rham des intersections compl\`etes, sur un corps de caract\'eristique quelconque, \`a celles de l'espace projectif.

On explicite ensuite quelques cons\'equences des r\'esultats obtenus. Au $n^{\circ}$ 2, on expose des r\'esultats qualitatifs d\'eduits de la formule num\'erique de Hirzebruch (2.2). Je tiens \`a remercier N. Katz de l'aide qu'il m'a apport\'ee dans la conception de ce num\'ero.

1.0. Soit $\underline{E}$ un module localement libre de rang $r+1$ sur un sch\'ema $S$. \underline{On suppose que} $r > 0$. Rappelons que $\mathbb{P}(\underline{E})$ d\'esigne l'espace projectif relatif sur $S$ des hyperplans de $\underline{E}$ . On d\'esignera par $p$ la projection de $\mathbb{P}(\underline{E})$ sur $S$ et par $\eta \in H^0(S,R^1 p_*(\Omega^1_{\mathbb{P}(\underline{E})/S}))$ la premi\`ere classe de Chern, style Hodge, du faisceau
%%%%%%%%%%  scan idx 47  |  folio 40  %%%%%%%%%%
inversible $\underline{O}(1)$. Rappelons que $\eta$ est l'image par la diff\'erentielle logarithmique $df/f : \underline{O}^* \longrightarrow \Omega^1_{\mathbb{P}(\underline{E})/S}$ de la classe de $\underline{O}(1)$ dans $H^0(S,R^1 p_* \underline{O}^*)$. Pour tout faisceau $\underline{F}$ de $\underline{O}$-modules sur $\mathbb{P}(\underline{E})$, on pose

\begin{equation*}
\underline{F}(n) = \underline{F} \otimes \underline{O}(1)^{\otimes n} \qquad (n \in \mathbb{Z}) \qquad .
\end{equation*}
\underline{Th\'eor\`eme} 1.1. \quad (i) \underline{Les faisceaux} $R^i p_*(\Omega^j_{\mathbb{P}(\underline{E})/S}(n))$ \underline{sont localement libres, et leur formation est compatible \`a tout changement de base.}

(ii) \underline{Pour} $0 \leq i \leq r$, \underline{la section} $\eta^i$ \underline{de} $R^i p_*(\Omega^i_{\mathbb{P}(\underline{E})/S})$ \underline{d\'efinit un isomorphisme entre ce faisceau et} $\underline{O}_S$ . \underline{De plus,}

\begin{equation*}
R^i p_* \ \Omega^j_{\mathbb{P}(\underline{E})/S} \ = \ 0 \qquad \underline{\mathrm{pour}}\ i \neq j \ \underline{\mathrm{ou}}\ i > r \qquad .
\end{equation*}
(iii) \underline{Si} $n \neq 0$, \underline{alors,} \underline{en dehors des deux cas suivants,} $R^i p_*(\Omega^j_{\mathbb{P}(\underline{E})/S}(n))$ \underline{est nul} :

\begin{equation}\tag{iii$_{a}$}
i = 0 \quad \underline{\mathrm{et}}\ n > j
\end{equation}
\begin{equation}\tag{iii$_{b}$}
i = r \quad \underline{\mathrm{et}}\ n < j-r \ .
\end{equation}
Soient $\mathbb{V}(\underline{E}) = \mathrm{Spec}\ (\mathrm{Sym}^*\ \underline{E})$ l'espace affine d\'efini par $\underline{E}$, $e$ sa section nulle, $\mathbb{V}(\underline{E})^*$ le compl\'ementaire $\mathbb{V}(\underline{E})^* = \mathbb{V}(\underline{E}) \setminus e(S)$, $a$ et $a'$ les projections de $\mathbb{V}(\underline{E})$ et $\mathbb{V}(\underline{E})^*$ sur $S$ et $q$ la projection de $\mathbb{V}(\underline{E})^*$ sur $\mathbb{P}(\underline{E})$ :

\begin{center}
\begin{tikzcd}[row sep=large, column sep=large]
\mathbb{V}(\underline{E}) \arrow[dr, "a"] & \mathbb{V}(\underline{E})^* \arrow[l, hook'] \arrow[r, "q"] \arrow[d, "a'"] & \mathbb{P}(\underline{E}) \arrow[dl, "p"] \\
& S \arrow[ul, bend right, "e"'] &
\end{tikzcd}
\end{center}
% STRUCTURE: 4 nodes, 6 arrows. NODES: N1 = V(E) [blackboard-bold V, E underlined], top-left; N2 = V(E)*
%   , top-centre; N3 = P(E) [blackboard-bold P, E underlined], top-right; N4 = S , bottom-centre, below
%   N2. ARROWS: (1) N2 -> N1, horizontal, pointing LEFT, arrowhead at N1, hook/curl at the tail end
%   (right, at N2) = hooked arrow (open immersion), UNLABELLED. (2) N2 -> N3, horizontal, pointing
%   RIGHT, plain arrowhead, label "q" set ABOVE the shaft. (3) N2 -> N4, vertical, pointing DOWN, plain
%   arrowhead at S, label "a'" set to the RIGHT of the shaft. (4) N1 -> N4, straight diagonal
%   down-right, plain arrowhead at S, label "a" set to the RIGHT/ABOVE the shaft. (5) N3 -> N4, straight
%   diagonal down-left, plain arrowhead just right of S, label "p" set BELOW/RIGHT of the shaft. (6) N4
%   -> N1, CURVED arrow bowing out to the left and downward, starting near S at bottom, sweeping left
%   then up, plain arrowhead at N1 pointing UP; label "e" set to the LEFT of the curve. No hooks, no
%   double-heads, no iso labels, no dashed lines anywhere else. A full stop (period) is printed at the
%   right margin on the S row, after the diagram; it is punctuation, not a node, so it is not
%   represented in the tikzcd.
L'espace affine \'epoint\'e $\mathbb{V}(\underline{O}(1))^*$ sur $\mathbb{P}(\underline{E})$ est canoniquement isomorphe \`a $\mathbb{V}(\underline{E})^*$, de sorte que, pour tout faisceau quasi-coh\'erent de $\underline{O}$-modules $\underline{F}$ sur $\mathbb{P}(\underline{E})$, on a canoniquement 


%%%%%%%%%%  scan idx 48  |  folio 41  %%%%%%%%%%

\begin{equation}\tag{1.1.1}
q_* \, q^* \, \underline{F} \;\simeq\; \sum_{n \in \mathbb{Z}} \underline{F}(n) \;\; .
\end{equation}
Cette graduation de $q_* \, q^* \, \underline{F}$ peut se déduire de l'action, sur ce faisceau, du groupe $\mathbb{G}_m$ des homothéties de $\mathbb{V}(\underline{E})$.

Les champs de vecteurs relatifs sur $\mathbb{V}(\underline{E})$ s'identifient aux $S$-morphismes de $\mathbb{V}(\underline{E})$ dans $\mathbb{V}(\underline{E})$ (identifié à l'espace tangent à l'origine). Le champ de vecteur $X$ défini par l'application identique de $\mathbb{V}(\underline{E})$ est invariant par homothéties, Ce champ de vecteur est tangent aux fibres de $q$ et définit une trivialisation

\begin{equation}\tag{1.1.2}
X \;:\; \Omega_{\mathbb{V}(\underline{E})^*/\,\mathbb{P}(\underline{E})} \longrightarrow \underline{O}_{\mathbb{V}(\underline{E})^*} \;\; .
\end{equation}
Si les $x_i$ forment une base de $\underline{E}$, on a

\begin{equation*}
X \;=\; \sum x_i \, dx^i \;\; .
\end{equation*}
Puisque $\Omega^1_{\mathbb{V}(\underline{E})^*/\mathbb{P}(\underline{E})}$ est de rang un, la suite exacte

\begin{equation*}
0 \longrightarrow q^*\Omega^1_{\mathbb{P}(\underline{E})/S} \longrightarrow \Omega^1_{\mathbb{V}(\underline{E})^*/S} \longrightarrow \Omega^1_{\mathbb{V}(\underline{E})^*/\mathbb{P}(\underline{E})} \longrightarrow 0
\end{equation*}
induit des suites exactes

\begin{equation*}
0 \longrightarrow q^*\Omega^i_{\mathbb{P}(\underline{E})/S} \longrightarrow \Omega^i_{\mathbb{V}(\underline{E})/S} \longrightarrow \Omega^1_{\mathbb{V}(\underline{E})^*/\mathbb{P}(\underline{E})} \otimes q^*\Omega^{i-1}_{\mathbb{P}(\underline{E})/S} \longrightarrow 0 \;\; .
\end{equation*}
Celles-ci se réécrivent (1.1.2)

\begin{equation}\tag{1.1.3}
0 \longrightarrow q^*\Omega^i_{\mathbb{P}(\underline{E})/S} \longrightarrow \Omega^i_{\mathbb{V}(\underline{E})^*/S} \xrightarrow{\;X\lfloor\;} q^*\Omega^{i-1}_{\mathbb{P}(\underline{E})/S} \longrightarrow 0 \;\; .
\end{equation}
On a canoniquement

\begin{equation}\tag{1.1.4}
\Omega^i_{\mathbb{V}(\underline{E})/S} \;=\; a^* \overset{i}{\wedge} \underline{E} \;\; .
\end{equation}
Les suites exactes courtes (1.1.3) s'obtiennent donc en fractionnant une suite exacte longue sur $\mathbb{V}(\underline{E})^*$

\begin{equation}\tag{1.1.5}
\ldots \longrightarrow {a'}^{*} \overset{i}{\wedge} \underline{E} \xrightarrow{\;X\lfloor\;} {a'}^{*} \overset{i-1}{\wedge} \underline{E} \xrightarrow{\;X\lfloor\;} {a'}^{*} \overset{i-2}{\wedge} \underline{E} \longrightarrow \ldots \;\; .
\end{equation}

%%%%%%%%%%  scan idx 49  |  folio 42  %%%%%%%%%%

Cette suite exacte longue n'est autre que la restriction à $\mathbb{V}(\underline{E})^*$ de la résolution de Kozul de $\underline{O}_{e(S)}$ sur $\mathbb{V}(\underline{E})$.

Projetons la suite exacte longue (1.1.5) sur $\mathbb{P}(\underline{E})$ par le morphisme affine $q$. En vertu de (1.1.1) et de la formule $q_*{a'}^{*} = q_*q^*p^*$, on obtient une suite exacte longue

\begin{equation}\tag{1.1.6}
\ldots \longrightarrow \sum_{n \in \mathbb{Z}} p^* \overset{i}{\wedge} \underline{E}(n) \xrightarrow{\;X\lfloor\;} \sum_{n \in \mathbb{Z}} p^* \overset{i-1}{\wedge} \underline{E}(n) \longrightarrow \ldots \;\; .
\end{equation}
Celle-ci se fractionne en suite exactes courtes, images de (1.1.3),

\begin{equation}\tag{1.1.7}
0 \longrightarrow \sum_{n \in \mathbb{Z}} \Omega^i_{\mathbb{P}(\underline{E})/S}(n) \longrightarrow \sum_{n \in \mathbb{Z}} p^* \overset{i}{\wedge} \underline{E}(n) \xrightarrow{\;X\lfloor\;} \sum_{n \in \mathbb{Z}} \Omega^{i-1}_{\mathbb{P}(\underline{E})/S}(n) \to 0 \;\; .
\end{equation}
Les suites exactes (1.1.7) sont sommes de suites exactes

\begin{equation}\tag{1.1.8}
0 \longrightarrow \Omega^i_{\mathbb{P}(\underline{E})/S}(n) \longrightarrow p^* \overset{i}{\wedge} \underline{E}(n-i) \xrightarrow{\;X\lfloor\;} \Omega^{i-1}_{\mathbb{P}(\underline{E})/S}(n) \longrightarrow 0 \;\; .
\end{equation}
Rappelons (EGA III 2.1.15 ou FAC) que les $R^i p_* \underline{O}(k)$ sont localement libres de formation compatible à tout changement de base, et que

\begin{equation}\tag{1.1.9}
R^i p_* \, \underline{O}(k) = 0 \ \text{si } i \neq 0,\ r, \ \text{ou si } i = 0,\ k < 0, \ \text{ou si } i = r,\ k > -r-1,
\end{equation}
et que $R^i p_* p^* \overset{j}{\wedge} \underline{E}(k)$ est donc nul sous ces mêmes conditions.

Les suites exactes (1.1.8) fournissent une résolution

\begin{equation}\tag{1.1.10}
0 \longrightarrow \Omega^i_{\mathbb{P}(\underline{E})/S}(n) \xrightarrow{\;E\;} p^* \overset{i}{\wedge} \underline{E}(n-i) \xrightarrow{\;X\lfloor\;} \ldots \longrightarrow p^* \overset{0}{\wedge} \underline{E}(n-0) \longrightarrow 0
\end{equation}
de $\Omega^i_{\mathbb{P}(\underline{E})/S}(n)$. La suite spectrale définie par cette résolution dégénère $(E_1 = E_\infty)$ pour des raisons de degré (1.1.9) ; pour $0 < j < r$, elle fournit

\begin{equation}\tag{1.1.11}
R^j p_* \Omega^i_{\mathbb{P}(\underline{E})/S}(n) = \mathcal{H} \, ( \ldots \xrightarrow{\;X\lfloor\;} p_* \overset{i-j}{\wedge} \underline{E}(n-i+j) \xrightarrow{\;X\lfloor\;} \ldots ) \;\; .
\end{equation}
Puisque $r \geq 1$, un argument de profondeur montre que

\begin{equation}\tag{1.1.12}
\sum_{n \in \mathbb{Z}} p_*p^* \overset{i}{\wedge} \underline{E}(n) = p_*q_*q^*p^* \overset{i}{\wedge} \underline{E} = {a'}_{*}{a'}^{*} \overset{i}{\wedge} \underline{E} = a_*a^* \overset{i}{\wedge} \underline{E} \;\; .
\end{equation}

%%%%%%%%%%  scan idx 50  |  folio 43  %%%%%%%%%%

 de sorte que $p_*p^* \wedge^i \underline{E}(n)$ est la partie homogène de poids $n + i$ de $a_*a^* \wedge^i \underline{E}$, soit $\mathrm{Sym}^n(\underline{E}) \otimes \wedge^i \underline{E}$ . Si on se rappelle que (1.1.5) provient de la résolution de Koszul, on trouve par (1.1.11) que, pour $0 < j < r$, $R^j p_* \Omega^i_{\mathbb{P}(\underline{E})/S}(n)$ n'est non nul que si $i-j = 0$ et $n-i+j = 0$, i.e. si $n = 0$ et $i = j$. Dans ce cas, de plus, $R^i p_* \Omega^i_{\mathbb{P}(\underline{E})/S}$ est isomorphe à $\underline{O}_S$ .

Pour $j = 0$, on déduit encore de (1.1.10) que

\begin{equation}\tag{1.1.13}
p_*\Omega^i_{\mathbb{P}(\underline{E})/S}(n) \;=\; \mathrm{Ker}(X \llcorner \;:\; p_* \wedge^i \underline{E}(n-i) \longrightarrow p_* \wedge^{i-1} \underline{E}(n-i+1)) \qquad ,
\end{equation}
et le même argument montre que $X \llcorner$ est injectif lorsque $n-i \leq 0$.

En vertu de(1.1.8) pour $n = 0$, $i = r+1$, on a $\Omega^r_{\mathbb{P}(\underline{E})/S} \simeq p^* \wedge^{r+1} \underline{E}(-r-1)$. Lorsque $S$ est le spectre d'un corps, les assertions de nullité de groupes de cohomologie de 1.1 (ii) et (iii) résultent de ce qui précède pour $j \neq r$, et s'en déduisent par dualité ([2] V 11.2 p. 211) pour $j = r$. En particulier, quels que soient $j$ et $n$, il existe au plus un $i$ tel que $R^i p_* \Omega^j_{\mathbb{P}(\underline{E})/S}(n) \neq 0$, ce qui implique (i) puisque $\Omega^j_{\mathbb{P}(\underline{E})/S}$ est plat sur $S$ (le plus dur est de trouver la référence dans EGA III ; je suggère : 6.10.5 + 7.4.1 + 7.5.5 + un passage à la limite).

Sachant que les $R^i p_* \Omega^j_{\mathbb{P}(\underline{E})/S}$ sont des faisceaux inversibles pour $0 \leq i \leq r$, il reste seulement à prouver que, pour $S$ le spectre d'un corps, on a $\eta^i \neq 0$. En effet $\eta^r$ est la classe de cohomologie d'un point rationnel (intersection de $r$ hyperplans), donc une base de $H^r(\mathbb{P}^r,\Omega^r)$ (FGA exp. 149 n$^{\circ}$4).

1.2. Il est possible de calculer, directement à partir de (1.1.10) ou (1.1.13), les caractéristiques d'Euler-Poincaré des $\Omega^i(n)$. Le résultat, peu appétissant, est

\begin{equation*}
\chi(\mathbb{P}^r_k,\Omega^j_{\mathbb{P}^r_k}(n)) \;=\; \sum_{a \leq i \leq j} \; (-1)^{i-j} \binom{r+1}{i}\binom{r+n-i}{r} \; .
\end{equation*}

%%%%%%%%%%  scan idx 51  |  folio 44  %%%%%%%%%%

\underline{Proposition} 1.3. \underline{Soient} $f : X \longrightarrow S$ \underline{un morphisme propre et lisse,} $\underline{W}$ \underline{un faisceau localement libre de rang} $c$ \underline{sur} $X$, $\underline{F}$ \underline{un faisceau cohérent sur} $X$, $d$ \underline{un entier} $> 0$ \underline{et} $s : \underline{W} \longrightarrow \underline{O}_X$ \underline{une section de} $\underline{W}^{\vee}$. \underline{On suppose que} $S$ \underline{est noethérien et que}

(a) \underline{Le sous-schéma} $H$ \underline{de} $X$ \underline{défini par l'annulation de} $s$ \underline{est lisse sur} $S$.

(b) \underline{Localement sur} $X$, \underline{si les} $s_i$ \underline{sont les coordonnées de} $S$ \underline{dans une base de} $\underline{W}^{\vee}$ \underline{les} $s_i$ \underline{forment une suite} $\underline{O}_X$\underline{-régulier et} $\underline{F}$\underline{-régulière.}

(c) \underline{Quelle que soit la suite d'entiers} $k_i$, \underline{non tous nuls, on a}

\begin{equation*}
R^i f_* \; (\underset{i}{\otimes} \; \wedge^{k_i} \; \underline{W} \otimes \Omega^j_{X/S} \otimes \underline{F}) \;=\; 0 \;\; \underline{\text{pour}} \;\; i+j < d \qquad .
\end{equation*}
\underline{Alors, on a} :

\begin{equation*}
R^i f_* \; (\Omega^j_{H/S} \otimes \underline{F}) \; \overset{\sim}{\longrightarrow} \; R^i f_*(\Omega^j_{X/S} \otimes \underline{F}) \;\; \underline{\text{pour}} \;\; i+j < d-c,
\end{equation*}
\begin{equation}\tag{ii}
R^i f_* \; (\Omega^j_{H/S} \otimes \underline{F}) \; \hookrightarrow \; R^i f_*(\Omega^j_{X/S} \otimes \underline{F}) \;\; \underline{\text{pour}} \;\; i+j = d-c \qquad .
\end{equation}
Le faisceau conormal $N_{H/X}$ de $H$ dans $X$ est isomorphe à la restriction de $\underline{W}$ à $H$, de sorte que la suite exacte

\begin{equation}\tag{1.3.1}
0 \longrightarrow N_{H/X} \longrightarrow \Omega^1_{X/S} \otimes \underline{O}_Y \longrightarrow \Omega^1_{H/S} \longrightarrow 0
\end{equation}
définit une filtration décroissante de $\Omega^j_{X/S} \otimes \underline{O}_Y$ dont les quotients successifs sont donnés par

\begin{equation}\tag{1.3.2}
\mathrm{Gr}^k(\Omega^j_{X/S} \otimes \underline{O}_Y) \qquad \simeq \qquad \wedge^k \; \underline{W} \otimes \Omega^{j-k}_{H/S} \qquad .
\end{equation}
Il résulte de (b) que l'on a encore

\begin{equation}\tag{1.3.3}
\mathrm{Gr}^k(\Omega^j_{X/S} \otimes \underline{O}_Y \otimes \underline{F}) \;=\; \wedge^k \; \underline{W} \otimes \Omega^{j-k}_{H/S} \otimes \underline{F} \qquad .
\end{equation}
On prouvera (i) et (ii) par récurrence sur $j$, simultanément pour tous les $\underline{F}$ vérifiant (b) et (c). Ces assertions sont vides pour $j < 0$ ;
%%%%%%%%%%  scan idx 52  |  folio 45  %%%%%%%%%%
supposons donc (i) et (ii) prouvés pour $j' < j$. L'hypothèse de récurrence s'applique aux faisceaux $\stackrel{k}{\wedge} \underline{W} \otimes \underline{F}$ ; elle fournit, par (c) :

\begin{equation}\tag{1.3.4}
R^i f_* ( \stackrel{k}{\wedge}\underline{W} \otimes \Omega^{j-k}_{H/S} \otimes \underline{F}) = 0 \quad \mbox{pour } k > 0 \mbox{ et } i+j < d-c+k \quad .
\end{equation}
La filtration considérée plus haut du $\Omega^j_{X/S} \otimes \underline{O}_Y \otimes \underline{F}$ définit une suite spectrale

\begin{equation*}
E_1^{p,i-p} \; = \; R^i f_*( \stackrel{p}{\wedge} \underline{W} \otimes \Omega^{j-p}_{H/S} \otimes \underline{F}) \Longrightarrow R^i f_*(\Omega^j_{X/S} \otimes \underline{O}_Y \otimes \underline{F}) \quad .
\end{equation*}
Par (1.3.4), on a $E_1^{p,i-p} = 0$ si $p > 0$ et $i+j \leq d-c$ d'où l'on tire que

\begin{equation}\tag{1.3.5}
\left\{ \begin{array}{l} R^i f_*(\Omega^j_{X/S} \otimes \underline{O}_Y \otimes \underline{F}) \; \stackrel{\sim}{\longrightarrow} \; R^i f_*(\Omega^j_{H/S} \otimes \underline{F}) \mbox{ pour } i+j < d-c \\[4ex] R^i f_*(\Omega^j_{X/S} \otimes \underline{O}_Y \otimes \underline{F}) \; \lhook\joinrel\longrightarrow \; R^i f_*(\Omega^j_{H/S} \otimes \underline{F}) \mbox{ pour } i+j = d-c \quad . \end{array} \right.
\end{equation}
D'après l'hypothèse (b), pour une différentielle convenable, les $\stackrel{k}{\wedge} \underline{W} \otimes \Omega^j_{X/S} \otimes \underline{F}$ forment une résolution à gauche (de Koszul) de $\Omega^j_{X/S} \otimes \underline{O}_Y \otimes \underline{F}$. Cette filtration définit une suite spectrale

\begin{equation}\tag{1.3.6}
E_1^{-p,q} \; = \; R^q f_*( \stackrel{p}{\wedge}\underline{W} \otimes \Omega^j_{X/S} \otimes \underline{F}) \Longrightarrow R^{q-p} f_*(\Omega^j_{X/S} \otimes \underline{O}_Y \otimes \underline{F}) \quad .
\end{equation}
D'après l'hypothèse (c), on a $E_1^{-p,q} = 0$ pour $p > 0$ et $q+j < d$, ou, trivialement, pour $p \notin [0,c]$. Dès lors,

\begin{equation}\tag{1.3.7}
\left\{ \begin{array}{l} R^i f_*(\Omega^j_{X/S} \otimes \underline{F}) \; \stackrel{\sim}{\longrightarrow} \; R^i f_*(\Omega^j_{H/S} \otimes \underline{F}) \mbox{ pour } i+j < d-c \\[4ex] R^i f_*(\Omega^j_{X/S} \otimes \underline{F}) \; \lhook\joinrel\longrightarrow \; R^i f_*(\Omega^j_{H/S} \otimes \underline{F}) \mbox{ pour } i+j = d-c \quad . \end{array} \right.
\end{equation}
Les assertions (i) et (ii) résultent maintenant de (1.3.5) et (1.3.7).


%%%%%%%%%%  scan idx 53  |  folio 46  %%%%%%%%%%

1.4. Soient $\underline{E}$ un faisceau de modules localement libres de rang $r+1$ sur un schéma $S$, $\mathbb{P}(\underline{E})$ le fibré projectif correspondant, $d$ un entier $\leq r$ et $\underline{a} = (a_1 \ldots a_d)$ une suite d'entiers. Une \underline{intersection complète} relative dans $\mathbb{P}(\underline{E})$, de multidegré $\underline{a}$, est un sous-schéma de $\mathbb{P}(\underline{E})$, fibre par fibre de codimension $d$, qui, localement sur $S$ pour la topologie étale, est intersection de $d$ hypersurfaces de degrés respectifs $a_1 \ldots a_d$.

En vertu de 1.1, les hypothèses (b) et (c) de 1.3 sont vérifiées pour $X/S = \mathbb{P}(\underline{E})/S$, $\underline{F} = \underline{O}_X$ , $\underline{W} = \sum\limits_i \underline{O}(-a_i)$ et pour $s$ une section de $\check{\underline{W}}$ définissant une intersection complète relative $Y \subset \mathbb{P}(\underline{E})$ de multidegré $\underline{a}$.

\underline{Théorème} 1.5. \underline{Sous les hypothèses} 1.4, \underline{soit} $n = r-d$ \underline{la dimension relative de} $Y$ \underline{et} $\eta \in H^o(S,R^1f_*\Omega^1_{Y/S})$ \underline{la première classe de Chern de} $\underline{O}(1)$

\begin{center}
\begin{tikzcd}[column sep=large, row sep=large]
Y \arrow[rr, hook] \arrow[dr, dash, "f"'] & & \mathbb{P}(\underline{E}) \arrow[dl, "g"] \\
 & S &
\end{tikzcd}
\end{center}
% STRUCTURE: 3 nodes, exactly as on the page. Node 1 (top left): Y. Node 2 (top right):
%   \mathbb{P}(\underline{E}) (double-struck P, E underlined). Node 3 (bottom centre, below and between
%   the other two): S. 3 edges. Edge A: Y --> \mathbb{P}(\underline{E}), horizontal, pointing right,
%   hooked/mono arrow (open hook at the Y end, arrowhead at the \mathbb{P}(\underline{E}) end),
%   unlabelled. Edge B: Y --- S, diagonal running down and to the right, labelled f placed to the LEFT
%   of the line at mid-height; the scan shows NO arrowhead at either end, so it is a plain undirected
%   line (dash). Edge C: \mathbb{P}(\underline{E}) --> S, diagonal running down and to the left,
%   arrowhead at the S end (lower-left), labelled g placed to the RIGHT of the line at mid-height. A
%   free-standing period (sentence full stop) sits on the baseline of S, far to the right of the figure.
\underline{Si} $Y$ \underline{est lisse sur} $S$, \underline{alors} :

(i) \underline{Les} $R^i f_* \Omega^j_{Y/S}$ \underline{sont localement libres}, \underline{et leur formation est compatible à tout changement de base}.

(ii) \underline{La suite spectrale} $E_1^{pq} = R^q f_* \Omega^p_{Y/S} \Longrightarrow R^{p+q} f_*(\Omega^*_{Y/S})$ \underline{reliant les cohomologie de Hodge et de De Rham dégénère} $(E_1 = E_\infty)$, \underline{son aboutissement est localement libre et sa formation est compatible à tout changement de base}.

(iii)$_a$ \underline{On a}

\begin{equation*}
R^i f_* \Omega^j_{X/S} = 0 \quad \underline{\mbox{si}} \; i \neq j \; \underline{\mbox{et}} \; i+j \neq n \; .
\end{equation*}

%%%%%%%%%%  scan idx 54  |  folio 47  %%%%%%%%%%

(iii)$_b$ $\eta^i$ \underline{est une base de} $R^i f_* \Omega^i_{Y/S}$ \underline{si} $2i < n$

(iii)$_c$ \underline{si} $n<2i\leq 2n$ \underline{et si} $\alpha_i$ \underline{est la base de} $R^i f_* \Omega^i_{Y/S}$ \underline{duale de la base} $\eta^{n-i}$ de $R^{n-i} f_* \Omega^{n-i}_{Y/S}$, \underline{alors}

\begin{equation*}
\eta^i \;=\; \left( \prod_1^d \; a_i \right) \alpha_i \qquad .
\end{equation*}
(iii)$_d$ \underline{Si} $n = 2m$, \underline{alors} $\eta^m \neq 0$ (\underline{si} $S \neq \emptyset$) \underline{et le quotient} $R^m f_* \Omega^m_{Y/S}/\underline{O}_S . \eta^m$ \underline{est localement libre}.

Il suffit de prouver 1.5 dans un cas ``universel'', ce qui permet de supposer $S$ lisse sur $\mathrm{Spec}(\mathbb{Z})$. Dans ce cas, l'assertion (ii) résulte de (i), car elle est vraie sur $\mathbb{C}$, donc sur tout corps de caractéristique $0$, d'après la théorie de Hodge. De plus, $S$ étant réduit, (i) résulte de (iii), supposé démontré pour $S$ le spectre d'un corps : en effet, d'après (iii), pour chaque $j$ les fonctions

\begin{equation*}
s \longrightarrow \dim_{k(s)} \; (H^i(Y_s,\Omega^j_{Ys})
\end{equation*}
sont constantes sur $S$, sauf peut-être pour $i = j$, donc sont toutes constantes puisque leur somme alternée l'est (EGA III 7.9.2) et on conclut par EGA III 6.10.5 + 7.6.9 + 7.7.5.

En vertu de (1.4) et (1.3), on a

\begin{equation}\tag{1.5.1}
\left\{
\begin{array}{ccl}
R^i f_* \Omega^j_{Y/S} & \stackrel{\sim}{\longrightarrow} & R^i g_* \Omega^j_{\mathbb{P}(\underline{E})/S} \qquad \mbox{pour } i+j < n \\[3ex]
R^i f_* \Omega^j_{Y/S} & \lhook\joinrel\longrightarrow & R^i g_* \Omega^j_{\mathbb{P}(\underline{E})/S} \qquad \mbox{pour } i+j = n \qquad .
\end{array}
\right.
\end{equation}
Si $S$ est le spectre d'un corps, les assertions (iii)$_a$ (iii)$_b$ (iii)$_d$ résultent de (1.5.1) et de la dualité de Serre ([2] V 11.2 p.211) La classe de
%%%%%%%%%%  scan idx 55  |  folio 48  %%%%%%%%%%
cohomologie $\eta^n$ est la classe de cohomologie d'un $0$-cycle de degré $\prod_1^d a_i$ (le degré de l'intersection complète), de sorte que

\begin{equation*}
\eta^n \;=\; \left( \prod_1^d \; a_i \right) \alpha_n \qquad .
\end{equation*}
Par définition, on a $\eta^{n-i} \alpha_i = \alpha_n$, et(iii)$_c$ en résulte.

Les assertions (i) et (ii) sont donc vraies, et les assertions (iii), dans le cas universel, résultent de (i) et de (iii) supposé prouvé pour $S$ le spectre d'un corps.

\underline{Théorème} 1.6. \underline{Gardons les hypothèses et notations de} 1.4, 1.5, \underline{supposons} $Y$ \underline{lisse,} \underline{soit} $\ell$ \underline{un nombre premier premier aux caractéristiques résiduelles de} $S$, \underline{et désignons par} $\eta_\ell \in H^0(S,R^2 f_*( \mathbb{Z}_\ell(1)))$ \underline{la première classe de Chern} $\ell$\underline{-adique de} $\underline{O}(1)$ (SGA 5 VII). \underline{Alors}

(i) $R^i f_* \mathbb{Z}_\ell = 0$ \underline{si} $i$ \underline{est impair et distinct de} $n$

(ii) \underline{Si} $2i < n$, \underline{alors} $R^{2i} f_* \mathbb{Z}_\ell(i)$ \underline{est isomorphe à} $\mathbb{Z}_\ell$, \underline{et admet} $\eta_\ell^i$ \underline{pour base}

(iii) \underline{Si} $n < 2i \leq 2n$, \underline{alors} $R^{2i} f_* \mathbb{Z}_\ell(i)$ \underline{est isomorphe à} $\mathbb{Z}_\ell$, \underline{et admet} $\eta_\ell^i / \prod_1^d a_i$ \underline{pour base.}

(iv) \underline{Le} $\mathbb{Z}_\ell$\underline{-faisceau} $R^n f_* \mathbb{Z}_\ell$ \underline{est constant tordu constructible sans torsion et,} \underline{si} $n$ \underline{est pair,} \underline{alors} $\eta_\ell^{n/2} \neq 0$ (\underline{si} $S \neq \emptyset$) \underline{et} $R^n f_* \mathbb{Z}_\ell(n/2)/ \mathbb{Z}_\ell . \eta_\ell^{n/2}$ \underline{est sans torsion.}

Les $\mathbb{Z}_\ell$-faisceaux $R^i f_* \mathbb{Z}_\ell$ sont constants tordus constructibles, et leur formation est compatible à tout changement de base, d'après SGA 4 XVI 2.2. Il suffit donc de prouver 1.6 pour $S$ le spectre d'un corps algébriquement clos. Une application itérée de SGA 4 XIV 3.3 montre alors, pour tout $m$, que la flèche canonique

\begin{equation*}
H^q_Y \; (\mathbb{P}(\underline{E}), \; \mathbb{Z}/\ell^m) \longrightarrow H^q(Y, \; \mathbb{Z}/\ell^m)
\end{equation*}

%%%%%%%%%%  scan idx 56  |  folio 49  %%%%%%%%%%
est bijective pour $q > r + d$ et surjective pour $q = r + d$. De là, et de la dualité de Poincaré (SGA 4 XVIII) on déduit (SGA 5 VII) que la flèche de restriction

\begin{equation*}
H^q(\ \mathbb{P}(\underline{E}), \mathbb{Z}_\ell) \longrightarrow H^q(Y, \mathbb{Z}_\ell)
\end{equation*}
est bijective pour $q < n$, injective et de conoyau sans torsion pour $q = n$. Ceci prouve (i) pour $i < n$, (ii) et (iv). On vérifie comme en 1.5 que $\eta_\ell^n$ est $\prod_1^d a_i$ fois la classe canonique, et (i) et (iii) s'obtiennent dès lors par dualité de Poincaré (SGA 4 XVIII).

1.7. Si $Y$ est une intersection complète de dimension au moins 3 dans un espace projectif $\mathbb{P}^r(k)$ sur un corps $k$, on sait que $\mathrm{Pic}(Y)$ est isomorphe à $\mathbb{Z}$, et est engendré par $\underline{O}(1)$ (SGA 2 XII 3.7). En dimension 2, et si $Y$ est lisse, on sait encore prouver :

\underline{Théorème} 1.8. \underline{Soit} $Y$ \underline{une surface lisse \ intersection complète dans l'espace projectif} $\mathbb{P}^r(k)$ \underline{sur un corps algébriquement clos} $k$. \underline{Alors,} $\underline{\mathrm{Pic}}^\tau(Y) = 0$, \underline{et la classe de} $\underline{O}(1)$ \underline{n'est pas divisible dans le groupe de Néron-Severi} $\mathrm{NS}(Y)$.

Puisque $H^1(Y,\underline{O}) = 0$ (1.5 (iii)$_a$), on sait que $\underline{\mathrm{Pic}}^0(Y) = 0$. La nullité de $\underline{\mathrm{Pic}}^\tau(Y)$ résulte alors du lemme bien connu suivant, du fait que $H^2(Y, \mathbb{Z}_\ell(1))$ est sans torsion (1.6 (iv)) et du fait que $H^0(Y,\Omega^1_Y) = 0$ (1.5 (iii)$_a$).

\underline{Lemme} 1.9. \underline{Soit} $Y$ \underline{une variété algébrique complète sur un corps algébriquement clos} $k$.

(i) \underline{Si $\ell$ est premier à la caractéristique de} $k$, \underline{alors la $\ell$-torsion de} $\mathrm{NS}(Y)$ \underline{s'identifie à la $\ell$-torsion de} $H^2(Y, \mathbb{Z}_\ell(1))$.
%%%%%%%%%%  scan idx 57  |  folio 50  %%%%%%%%%%
(ii) \underline{Si} $k$ \underline{est de caractéristique non nulle} $p$ \underline{et si} $Y$ \underline{est lisse, alors le noyau de la multiplication par} $p$ \underline{dans} $\mathrm{Pic}(Y)$ \underline{s'identifie au sous-groupe de} $H^0(Y,\Omega^1_Y)$ \underline{formé des différentielles localement logarithmiques.}

Les suites exactes de Kummer

\begin{equation*}
0 \longrightarrow \mu_{\ell^n} \longrightarrow G_m \xrightarrow{\ x \mapsto x^{\ell^n}\ } G_m \longrightarrow 0
\end{equation*}
définissent des suites exactes

\begin{equation*}
0 \longrightarrow \mathrm{Pic}(Y)/\ell^n\ \mathrm{Pic}(Y) \longrightarrow H^2(Y,\mu_{\ell^n}) \longrightarrow (H^2(G_m))_{\ell^n} \longrightarrow 0 \quad ,
\end{equation*}
où $\mathrm{Pic}(Y)/\ell^n\ \mathrm{Pic}(Y) = \mathrm{NS}(Y) \otimes \mathbb{Z}/\ell^n$. Par passage à la limite projective, on trouve

\begin{equation*}
0 \longrightarrow \mathrm{NS}(Y) \otimes \mathbb{Z}_\ell \longrightarrow H^2(Y, \mathbb{Z}_\ell(1)) \longrightarrow \mathrm{Hom}(\mathbb{Q}_\ell / \mathbb{Z}_\ell,\ H^2(G_m)) \longrightarrow 0 \quad ,
\end{equation*}
et le conoyau est sans torsion puisque $\mathbb{Q}_\ell / \mathbb{Z}_\ell$ est divisible. L'assertion (i) en résulte.

Pour prouver (ii), il suffit de noter que la suite exacte

\begin{equation*}
0 \longrightarrow \underline{O}^*_Y \xrightarrow{\ x \longmapsto x^p\ } \underline{O}^*_Y \xrightarrow{\ df/f\ } \Omega^1_Y
\end{equation*}
définit une suite exacte courte

\begin{equation*}
0 \longrightarrow \underline{O}^*_Y \longrightarrow \underline{O}^*_Y \longrightarrow \Omega^1_{Y\,\mathrm{loc}\,\mathrm{log}} \longrightarrow 0
\end{equation*}
et de prendre la suite exacte longue de cohomologie correspondante. Ce raisonnement est dû à Cartier, qui prouve de plus que, si l'on désigne par $C$ l'opération de Cartier sur $H^0(Y,\Omega^1_Y)$, alors les différentielles localement logarithmique forment le noyau de $C - \mathrm{Id}$.

Supposons par l'absurde que la classe de $\underline{O}(1)$ dans $\mathrm{NS}(Y)$ soit de la forme $nu$, avec $n > 1$. Si $\ell$ divise $n$ et est premier à la caractéristique,
%%%%%%%%%%  scan idx 58  |  folio 51  %%%%%%%%%%
on a encore, en cohomologie $\ell$-adique

\begin{equation*}
\eta_\ell = n\ c_1(u)
\end{equation*}
ce qui contredit 1.6 (iv).

Enfin, si la caractéristique $p$ de $k$ est différente de $0$, et si $p$ divise $n$, on a encore, en cohomologie de Hodge

\begin{equation*}
\eta = n\ c_1(u) = 0 \qquad ,
\end{equation*}
ce qui contredit 1.5 (iii)$_d$ et achève la démonstration.


2. \underline{Résultats numériques}

2.1. Soient $\underline{a}=(a_i)_{1\leq i\leq d}$ une suite d'entiers $\geq 1$, et désignons par $V_n(\underline{a})$ une intersection complète lisse de dimension $n$ et de multidegré $\underline{a}$ dans un espace projectif $\mathbb{P}^{n+d}$ sur un corps $k$. On posera

\begin{equation*}
h^{pq}\ (\underline{a}) = \dim_k H^q(V_{p+q}(\underline{a}),\ \Omega^p)
\end{equation*}
et

\begin{equation*}
h_o^{pq}(\underline{a}) = h^{pq}(\underline{a}) - \delta_{p,q} \qquad .
\end{equation*}
Les $h_o^{pq}(\underline{a})$ sont donc les nombres de Hodge de la cohomologie primitive de dimension moitié d'une intersection complète de multidegré $\underline{a}$. Ils ne dépendent que de $\underline{a}$, $p$ et $q$ (1.5). On posera

\begin{equation}\tag{2.1.1}
H(\underline{a}) = \sum_{p,q\geq 0} h_o^{p,q}\ y^p z^q \quad \in \mathbb{Z}[[y,z]] \qquad .
\end{equation}
Pour $d = 1$ et $\underline{a} = (a)$, on pose $h^{pq}(a) = h^{pq}(\underline{a})$ et $H(a) = H(\underline{a})$. Pour $d = 0$, on pose $\sup(a_i) = 1$.

La formule suivante est le cas particulier $k = 0$ de la formule (2) p.160 de Hirzebruch [3] 


%%%%%%%%%%  scan idx 59  |  folio 52  %%%%%%%%%%

\begin{equation}\tag{2.2}
\sum_{n=0}^{\infty}\ \sum\ \chi(V_n(\underline{a}),\ \Omega^p)\ y^p\ z^{n+d}\ =\ \frac{1}{(1+yz)(1-z)}\ \prod_{1}^{d}\ \frac{(1+yz)^{a_i}-(1-z)^{a_i}}{(i+yz)^{a_i}+y(1-z)^{a_i}}
\end{equation}
Nous allons transformer cette formule en la suivante :

\underline{Théorème} 2.3. (Hirzebruch). \underline{Avec les notations} 2.1, \underline{on a}

\begin{equation*}
H(\underline{a}) = \frac{1}{(1+y)(1+z)}\left[\ \prod_{1\leq i\leq d}\ \frac{(1+y)^{a_i}\ -\ (1+z)^{a_i}}{-(1+y)^{a_i}z\ +\ (1+z)^{a_i}y}\ \ -1\ \right] \qquad .
\end{equation*}
En vertu de 1.5, on a

\begin{equation*}
\chi(V_n(\underline{a}),\Omega^p) = (-1)^p + (-1)^{n-p}\ h_o^{p,n-p}(\underline{a})
\end{equation*}
de sorte que le premier membre $I$ de (2.2) se réécrit

\begin{equation*}
\begin{aligned}
I &= \sum_{n=0}^{\infty}\ \sum_{p+q=n} ((-1)^p+(-1)^q\ h_o^{p,q}(\underline{a}))\ y^p z^{p+q+d}\\
&= \sum_{p,q\geq 0} (-1)^p(yz)^p z^q z^d\ +\ \sum_{p,q\geq 0} (-1)^q\ h_o^{p,q}(a)\ (yz)^p z^q\ z^d\\
&= z^d\left[\frac{1}{(1+yz)(1-z)}\ +\ H(\underline{a})(yz,-z)\right] \qquad .
\end{aligned}
\end{equation*}
Le deuxième membre $II$ de (2.2) se réécrit

\begin{equation*}
II = \frac{z^d}{(1+yz)(1-z)}\ \prod_{1}^{d}\ \frac{(1+yz)^{a_i}\ -\ (1-z)^{a_i}}{(1+yz)^{a_i}z+\ yz(1-z)^{a_i}}
\end{equation*}
de sorte que, simplifiant par $z^d$ dans $I = II$, on trouve que

\begin{equation*}
H(\underline{a})(yz,-z) = \frac{1}{(1+yz)(1-z)}\left[\ \prod_{1}^{d}\ \frac{(1+yz)^{a_i}\ -\ (1-z)^{a_i}}{(1+yz)^{a_i}z\ +\ (1-z)^{a_i}yz}\ \ -1\right]
\end{equation*}
et 2.3 se déduit de cette formule par changement de variable.


%%%%%%%%%%  scan idx 60  |  folio 53  %%%%%%%%%%

\underline{Corollaire} 2.4. \underline{Avec les notations} 2.1,

(i) \quad \underline{pour} $d = 1$, \underline{on a}

\begin{equation*}
\begin{aligned}
H(a) \;&=\; -\ \frac{(1+y)^{a-1} - (1+z)^{a-1}}{(1+y)^{a}z \,-\, (1+z)^{a}y}\\[2ex]
&=\; \frac{\displaystyle\sum_{i,j\geq 0}\binom{a-1}{i+j+1}\ y^{i}z^{j}}{1 \,-\, \displaystyle\sum_{i,j\geq 1}\binom{a}{i+j}y^{i}\ z^{j}}\ .
\end{aligned}
\end{equation*}
(ii) \quad \underline{on a}, $|P|$ \underline{désignent le nombre d'élément d'une partie} $P$ \underline{de} $[1,d]$,

\begin{multline*}
H(\underline{a}) \;=\; \sum_{\substack{P \subset [1,d] \\ P \neq \emptyset}} (1+y)^{|P|-1}\ (1+z)^{|P|-1}\ \prod_{i \in P} H(a_{i})\ .
\end{multline*}
D'après (2.3), on a

\begin{equation*}
\begin{aligned}
H(a) \;&=\; \frac{-1}{(1+y)(1+z)}\left[\ \frac{(1+y)^{a} \,-\, (1+z)^{a}}{(1+z)^{a}\cdot z \,-\, (1+z)^{a}y} \;+\; 1\ \right]\\[2ex]
&=\; \frac{-1}{(1+y)(1+z)}\left[\ \frac{(1+y)^{a}(1+z)-(1+z)^{a}(1+y)}{(1+y)^{a}z \,-\, (1+z)^{a}y}\ \right]\\[2ex]
&=\; -\ \frac{(1+y)^{a-1} \,-\, (1+z)^{a-1}}{(1+y)^{a}z \,-\, (1+z)^{a}y}\ .
\end{aligned}
\end{equation*}
Soient $N$ et $D$ les numérateurs et dénominateurs de cette formule.

On a

\begin{equation*}
H(a) \;=\; -\,\frac{N}{D}\ , 
\end{equation*}

%%%%%%%%%%  scan idx 61  |  folio 54  %%%%%%%%%%

\begin{equation*}
 N \;=\; \sum \binom{a-1}{k}\ (y^{k} \,-\, z^{k}) \;=\; (y-z) \sum_{i,j\geq 0}\binom{a-1}{i+j+1}\ y^{i}\ z^{j}
\end{equation*}
\begin{equation*}
\begin{aligned}
D \;&=\; \sum \binom{a}{k}\ (y^{k}z \,-\, yz^{k}) \;=\; (z-y) \;+\; yz \sum_{k\geq 1}\binom{a}{k}\!\left(y^{k-1} \,-\, z^{k-1}\right)\\[2ex]
&=\; (y-z)\ (-1 \;+\; yz \sum_{i,j\geq 0}\binom{a}{i+j+2}\ y^{i}z^{j}\\[2ex]
&=\; (y-z)\ (-1 \;+\; \sum_{i,j\geq 1}\binom{a}{i+j}y^{i}\ z^{j})
\end{aligned}
\end{equation*}
et l'assertion (i) en résulte.

Inversons la formule de 2.3, pour $d = 1$. On trouve

\begin{equation*}
-\ \frac{(1+y)^{a} \,-\, (1+z)^{a}}{(1+y)^{a}z \,-\, (1+z)^{a}y} \;=\; 1 \;+\; (1+x)(1+y)\ H(a)\ .
\end{equation*}
Si l'on substitue cette formule dans 2.3, on obtient

\begin{equation*}
\begin{aligned}
H(\underline{a}) \;&=\; \frac{1}{(1+y)(1+z)}\left[\ \prod_{1}^{d}\ (1+(1+x)(1+y)\ H(a_{i}))\qquad -\ 1\ \right]\\[2ex]
&=\;\frac{1}{(1+y)(1+z)}\left[\ \sum_{\substack{P \subset [i,d] \\ P \neq \emptyset}}\ (1+x)^{|P|}(1+y)^{|P|}\ \prod_{i \in P} H(a_{i})\ \right]
\end{aligned}
\end{equation*}
d'où résulte (ii).

\underline{Théorème} 2.5. \underline{Avec les notations} 2.1

(i) \quad \underline{Si} $p+q = p'+q'$ \underline{et si} $p \leq p' \leq q' \leq q$, \underline{alors}

\begin{equation*}
h_{o}^{pq}\ (\underline{a}) \;\leq\; h_{o}^{p'q'}(\underline{a})
\end{equation*}
(ii) \quad \underline{Si} $p \leq q$, \underline{pour que} $h_{o}^{pq}\ (\underline{a}) \neq 0$, \underline{il faut et il suffit que} 


%%%%%%%%%%  scan idx 62  |  folio 55  %%%%%%%%%%

\begin{equation*}
 p \;>\; \left[\frac{p+q+d \;-\; \sum a_i}{\sup(a_i)}\right] \qquad .
\end{equation*}
\underline{où} [ ] \underline{désigne la partie entière du quotient.}

Soit (*) la propriété suivante d'une série formelle $P = \sum a_{pq} \, y^p \, z^q$ :

(*)  les $a_{pq}$ sont positifs, $a_{pq} = a_{qp}$, et si $p+q = p'+q'$, et $p \leqslant p' \leqslant q' \leqslant q$, alors $a_{pq} \leqslant a_{p'q'}$ .

L'assertion (i) résulte aussitôt des formules (2.4), de l'identité $(1-x)^{-1} = \sum_{n \geqslant 0} x^n$ , et du lemme suivant, dont la démonstration est laissée au lecteur :

\underline{Lemme} 2.5.1. \underline{L'ensemble des séries formelles vérifiant} (*) \underline{est stable par addition et multiplication.}

On prouvera (ii) tout d'abord dans le cas particulier des hypersurfaces ($d = 1$). La formule à démontrer se réécrit alors, pour $p \leqslant q$,

\begin{equation}\tag{2.5.2}
h_o^{pq}(a) \neq 0 \quad \Longleftrightarrow \quad p \geqslant \left[\frac{p+q+1}{a}\right]
\end{equation}
\begin{equation*}
(\; \Longleftrightarrow \quad (p+1) \; a \geqslant p+q+2 \quad \mbox{trivialement}).
\end{equation*}
De la deuxième formule 2.4 (i), on déduit que $h_o^{pq}$ est non nul si et seulement si il existe des couples $(e,f)$ et $(c_i,d_i)$ tels que

\begin{equation}\tag{2.5.3}
\left\{
\begin{array}{l}
(e,f) \geqslant (0,0) \quad \mbox{et } e+f \leqslant a-2 \\[3mm]
(c_i,d_i) \geqslant (1,1) \mbox{ et } c_i+d_i \leqslant a \\[3mm]
(p,q) = (e,f) + \sum (c_i,d_i) \qquad .
\end{array}
\right.
\end{equation}

%%%%%%%%%%  scan idx 63  |  folio 56  %%%%%%%%%%

Pour $p+q$ donné, le plus petit $p$ possible s'obtient pour $e = 0$ et $c_i = 1$, et est le nombre d'indices $i$. Ce $p$ est donc la borne inférieure des $p$ pour lesquels on a

\begin{equation*}
p \;\geqslant\; \frac{(p+q)-(a-2)}{a} \qquad , \; \mbox{i.e.} \quad (p+1)a \geqslant p+q+2 \qquad .
\end{equation*}
Passons au cas général. Si $\underline{a}'$ se déduit de $\underline{a}$ en supprimant les $a_i$ égaux à 1, alors $H(\underline{a}) = H(\underline{a}')$. Vu la forme de la formule 2.5 (ii), on peut supposer que pour tout $i$, $a_i > 1$. Si $d = 0$, alors $h_o^{pq} = 0$ et l'assertion est correcte; on ne restreint pas la généralité en supposant que $d > 0$ et que $a_d = \sup(a_i)$.

En vertu de 2.4 (ii), et (2.5.2), pour que $h_o^{pq}(\underline{a}) \neq 0$, il faut et il suffit qu'il existe une partie non vide $P$ de $[1,d]$, un couple $(e,f)$ et des couples $(c_i,d_i)_{i \leqslant i \leqslant d}$ tels que

\begin{equation}\tag{2.5.4}
\left\{
\begin{array}{l}
(|P|-1, \; |P|-1) \geqslant (e,f) \geqslant (0,0) \; , \\[3mm]
c_i,d_i \geqslant \left[{\displaystyle \frac{c_i+d_i+1}{a_i}}\right] \quad ; \quad (c_i,d_i) \geqslant (0,0) \quad , \\[3mm]
(p,q) = (e,f) + \sum (c_i,d_i) \qquad .
\end{array}
\right.
\end{equation}
Pour $n = p+q$ fixé, cherchons la plus petite valeur de $p$ possible. Si $n_i = c_i + d_i$, pour un tel $p$, les conditions (2.5.4) se réécrivent, en terme des $c_i$, $n_i$, $e$ et $f$ 


%%%%%%%%%%  scan idx 64  |  folio 57  %%%%%%%%%%

\begin{equation}\tag{2.5.5}
\left\{
\begin{array}{l}
a_i = 2 \Longrightarrow n_i \mbox{ pair} \qquad \mbox{(sinon (2.5.4) est impossible)} \\[3mm]
c_i \geq \left[ \displaystyle\frac{n_i+1}{a_i} \right] \quad ; \quad c_i \leq n_i \quad , \\[3mm]
(|P|-1, \; |P|-1) \geq (e,f) \geq (0,0) \quad , \\[3mm]
n = e+f+ \sum n_i \quad , \\[3mm]
p = e + \sum c_i \quad .
\end{array}
\right.
\end{equation}
Cette description montre tout d'abord que le plus petit $p$ est obtenu pour $P = [1,d]$, cas auquel on se limitera. Si on a

\begin{equation}\tag{2.5.6}
\sum (a_i-2) + (d-1) \geq n,
\end{equation}
alors il existe une solution avec $p = c_i = e = 0$. Sinon, si $a_d \neq 2$, le mieuxqu'on puisse faire est de prendre $n_i = a_i - 2$ pour $i \neq d$, de sorte que $c_i = 0$ pour $i \neq d$, $e = 0$, $f = d-1$ et $n_d = n - (d-1) - \sum\limits_{i \neq d} (a_i-2)$. On a alors

\begin{equation}\tag{2.5.7}
p = \left[ \displaystyle\frac{n-(d-1) - \sum\limits_{i \neq d} (a_i-2) + 1}{a_d} \right] \qquad ,
\end{equation}
et donc, pour $p \leq q$, on a $h^{pq}_{o}(\underline{a}) \neq 0$ si et seulement si

\begin{equation}\tag{2.5.8}
p \geq \left[ \displaystyle\frac{n- \sum\limits_{i \neq d} (a_i-1) + 1}{a_d} \right] \qquad .
\end{equation}
La formule (2.5.8) se réduit à (2.5.2) si $d = 1$ et est trivialement vraie si (2.5.6) est vérifié. Si $d \geq 2$, si (2.5.6) n'est pas vérifié, et si tous les $a_i$ sont égaux à $2$, le mieux qu'on puisse faire est encore de
%%%%%%%%%%  scan idx 65  |  folio 58  %%%%%%%%%%
prendre $n_i = 0$ $(i \neq d)$, $e = 0$, $f = d-1$ ou $d-2$, congru à $n$ modulo $2$, et $n_d = n-f$. Si $n \equiv d-1$ $(2)$, on prouve (2.5.7) comme plus haut. Sinon, on note que

\begin{equation*}
\left[ \displaystyle\frac{n - (d-2) + 1}{2} \right] = \left[ \displaystyle\frac{n - (d-1) + 1}{2} \right]
\end{equation*}
de sorte que (2.5.7), et par là (2.5.8), est valable même si $a_d = 2$.

Reste à vérifier que (2.5.8) équivaut à la formule 2.5 (ii). En effet, on a

\begin{equation*}
\left[ \displaystyle\frac{n - \sum\limits_{i \neq d} (a_i-1) + 1}{a_d} \right] = \left[ \displaystyle\frac{n + d - \sum\limits_{i \neq d} a_i}{a_d} \right] = \left[ \displaystyle\frac{p + q + d - \sum a_i}{a_d} \right] + 1 \quad .
\end{equation*}
2.6. \quad Si $X$ est une intersection complète de dimension $n$ dans $\mathbb{P}^r$, il résulte aussitôt de (2.5.8) que la partie primitive de $H^n(X)$ n'est nulle que si $X$ est une quadrique de dimension impaire.

\underline{Définition} 2.7. \underline{Si} $X$ \underline{est une variété algébrique propre et lisse sur} $\mathbb{C}$, \underline{le niveau de Hodge de} $H^n(X)$ \underline{est la borne supérieure des nombres} $q-p$ \underline{pour} $p+q = n$ \underline{et} $h^{pq}(X) \neq 0$.

Le niveau de Hodge $\ell$ est donc congru à $n$ modulo $2$, et est au plus égal à $\inf(n,\dim(X) - n)$. Il est $\geq 0$, sauf si $H^n(X) = 0$, auquel cas il est égal à $- \infty$.

\underline{Corollaire} 2.8. \underline{Reprenons les notations} 2.1, \underline{en faisant} $k = \mathbb{C}$, \underline{et en supposant que} $a_d = \sup (a_i)$. \underline{Soit} $\ell$ \underline{un entier} $0 \leq \ell < n$, \underline{congru} à $n$ \underline{modulo} $2$. \underline{Pour que le niveau de Hodge de} $H^n(V_n(\underline{a}))$ \underline{soit} $\leq \ell$ , \underline{il faut et il suffit que} 


%%%%%%%%%%  scan idx 66  |  folio 59  %%%%%%%%%%

\begin{equation*}
 \sum_{i \neq d} \, (a_i - 1) \;\leq\; (\ell + 1) \, - \, (a_d - 2) \, \frac{n - \ell}{2} \qquad .
\end{equation*}
En vertu de (2.5.8), si le niveau de Hodge est $\geq 0$, il est la borne supérieure de 0 et de

\begin{equation*}
n \, - \, 2 \left[ \frac{n \, - \, \sum_{i \neq d} \, (a_i - 1) \, + \, 1}{a_d} \right]
\end{equation*}
et l'inéquation

\begin{equation*}
n \, - \, 2 \left[ \frac{n \, - \, \sum_{i \neq d} \, (a_i - 1) \, + \, 1}{a_d} \right] \;\leq\; \ell
\end{equation*}
peut se réécrire

\begin{equation*}
\frac{n \, - \, \ell}{2} \;\leq\; \left[ \frac{n - \, \sum_{i \neq d} \, (a_i - 1) \; + \; 1}{a_d} \right] \qquad ,
\end{equation*}
ou encore

\begin{equation*}
\sum_{i \neq d} \, (a_i - 1) \;\leq\; 1 \, + \, n \, - \, a_d \; . \; \frac{n \, - \, \ell}{2} \;=\; \ell \, + \, 1 \, - \, (a_d - 2) \; \frac{n \, - \, \ell}{2} \qquad .
\end{equation*}
2.9. \quad A l'aide de la formule 2.8, il est facile d'énumérer, pour $\ell$ petit, les multidegrés des intersections complètes dans $\mathbb{P}^r$, de dimension $n > 1$, dont le $H^n$ est de niveau $\ell$. La table donnée dans la première version de cet exposé était néanmoins incorrecte, et celle donnée ci-dessous est extraite de celle plus complète donnée par M. Rapoport dans son article : Complément à l'article de P. Deligne "La conjecture de Weil pour les surfaces K 3", Inventiones Math. \underline{15} 227-236 (1972). Pour $\ell \leq 1$, les intersections complètes $V_n(\underline{a})$ de dimension $n$ et de multidegré $\underline{a} = (a_1,\ldots,a_d)$, avec $n > \ell$ , $d \geq 1$ et $a_i \geq 2$, de niveau de Hodge $\ell$ , sont les suivantes :


%%%%%%%%%%  scan idx 67  |  folio 60  %%%%%%%%%%

\begin{equation*}
\begin{array}{ll}
\ell \;=\; -\infty & V_{2k+1}(2) \qquad\qquad\qquad (k \geq 0) \\[2ex]
\ell \;=\; 0 & V_{2k}(2) \; (k \geq 1), \; V_{2k}(2,2) \; (k \geq 1), \; V_2(3) \\[2ex]
\ell \;=\; 1 & V_{2k+1}(2,2) \; (k \geq 1), \; V_{2k+1}(2,2,2) \; (k \geq 1) \\
& V_3(3), \; V_5(3), \; V_3(3,2), \; V_3(4) \; .
\end{array}
\end{equation*}
2.10. \quad Pour $\underline{a}$ petit, les fonctions génératrices des $h_o^{pq}(\underline{a})$ s'écrivent, d'après 2.4

\begin{equation*}
\begin{array}{rcl}
H(2) & = & \dfrac{1}{1-yz} \\[3ex]
H(3) & = & \dfrac{2 \, + \, y \, + \, z}{1 \, - \, 3yz \, - \, y^2 z \, - \, yz^2} \\[3ex]
H(2,2) & = & \dfrac{y \, + \, z}{(1-yz)^2} \; + \; \dfrac{2}{(1-yz)^2} + \; \dfrac{1}{1-yz}
\end{array}
\end{equation*}
de sorte que

\begin{equation*}
\begin{array}{rcll}
h_o^{pp}(2) & = & 1 & , \\[2ex]
h_o^{11}(3) & = & 6 & , \\[2ex]
h_o^{pp}(2,2) & = & 2p \, + \, 3 & .
\end{array}
\end{equation*}

%%%%%%%%%%  scan idx 68  |  folio 61  %%%%%%%%%%


\underline{BIBLIOGRAPHIE}

[1] Y. AKIZUKI and S. NAKANO. Note on Kodaira - Spencer's proof of Lefschetz theorems - Proc. Jap. Acad. 30, 266-272, (1954).

[2] R. HARTSHORNE. Residues and duality. Lecture notes 20 Springer 66.

[3] F. HIRZEBRUCH. Topological methods in algebraic geometry 3é ed. Springer-Verlag 1966.

[4] K. KODAIRA. On a differential geometric method in the theory of analytic stacks. Proc. Nat. Acad. Sci. USA 39, 1268-1273, (1953).
